
\documentclass[a4paper]{article}
\usepackage[utf8]{inputenc}
\usepackage[T1]{fontenc}
\usepackage[french]{babel}
\usepackage{fixltx2e}
\usepackage{nth}
\usepackage{listings}
\usepackage{tikz}
\usepackage{fancyhdr}
\usepackage[margin=1.0in]{geometry}
\usepackage{hyperref}

\pagestyle{fancy}
\lhead{{\sc{OCaml}}\\ {\sc TP-Spé} Traitement d'image - 11 oct. 2013}
\rhead{{\small \sc{GConfs}}\\ {\sc Epita}}
\rfoot{\includegraphics[width=0.15\linewidth]{gconfs.png}}

 \begin{document}
\begin{center}
  {\Large {\bf Le traitement d'image en
  \textsc{OCaml}}}
\end{center}

\bigskip

\section*{Introduction}

Ce soir (voir nuit pour les plus téméraires) nous allons nous intéresser au traitement d'images en \textsc{OCaml}, et c'est pour cela que 
vous êtes restés me direz-vous !\\
Tout d'abord laissez moi vous dire que durant ce TP nous allons principalement vous introduire les
 notions nécéssaire à faire un jeu avec une interface en \textsc{OCaml} avec OcamlSdl.
Le TP sera découpé en plusieurs étapes/paliers afin que vous puissiez réaliser, dans un premier temps un Pong, puis pour ceux qui avance et qui restent un jeu de plate-forme\\\\
Durant ce tp 2 armes seulement vous seront nécéssaires:
\begin{itemize}
\item Les bibliothèques SDL (normalement déjà installés sur vos racks)
\item La documentation de SDL (\url{ocamlsdl.sourceforce.net} partie Ressouces, puis Ocamldoc)
\end{itemize}

\section{Premiers pas avec SDL}
	\subsection{Le chargement des images}
		\begin{verbatim}
		val load_image : string -> Sdlvideo.surface
		\end{verbatim}
	\subsection{L'affichage des images et les DOUBLEBUF}
		\begin{verbatim}
		val set_video_mode : w:int -> h:int -> video_flag list -> surface
		val flip : surface -> unit
		\end{verbatim}

\section{Organisation d'une classe d'affichage}
\subsection{Utilisation de la fenêtre}
\subsection{Methodes de modification de la fenêtre}

\section{Boucle de rafraîchissement}

\section{Les récuperations des entrées clavier}
	\subsection{Récuperation des touches de déplacement}
	Sdlkey
	\begin{verbatim}
	val get_key_state : unit-> (int, Bigarray.int8_unsigned_elt, Bigarray.c_layout) Bigarray.Array1.t
	val int_of_key
	\end{verbatim}

	\subsection{Déplacement des raquettes en fonction des touches}

\section{Mise en place du déplacement d'une balle}
	\subsection{Deplacement aléatoire}
	\subsection{Rebond sur murs et raquettes}
	\subsection{Zones de perte}

\end{document}
